
%\needspace{4\baselineskip}
\color{uniblue}{\subsection{Частотное автоматическое повторное включение (ЧАПВ)}}
\color{black}

\begin{enumerate}

\item
Частотное автоматическое повторное включение используется для восстановления электроснабжения потребителей электроэнергии, которые были отключены во время срабатывания АЧР. В устройстве частотное автоматическое повторное включение представлено функцией «ЧАПВ».

\item
Основные характеристики функции:
\begin{itemize}
    \item погрешность ИО частоты не более 20 мГц;
    \item время срабатывания и возврата ИО частоты не более 60 мс.
\end{itemize}

\item
Готовность ЧАПВ реализуется при наличии включенного положения выключателя «В: Включен», отсутствие блокировки включения из логики КСВ или внешнего запрета ЧАПВ. При наличии данных условий запускается независимая выдержка времени таймера Т1 и происходит срабатывание триггера готовности ЧАПВ.

\item
В функции предусмотрено два режима работы, определяемые уставкой SGF2:
\begin{itemize}
    \item по входу;
    \item измеренный.
\end{itemize}

\item
В режиме «По входу» если уставка SGF3 выставлена в положение тип пуска «Внутренний» то пусковые условия будут формироваться при воздействии аварийного отключения выключателя от функции АЧР. Действие отключения от АЧР продлевается на время таймера Т3 равного 1 с, тем самым приводя к срабатыванию RS-триггера. По окончанию действия функции АЧР собираются условия для разрешения работы ЧАПВ.

\item
В режиме «По входу», при установке SGF3 в положение типа пуска «Внешний», условия для пуска формируются таким же образом, как описано ранее. Однако основным требованием остается наличие внешнего дискретного сигнала «ЧАПВ».

\item
В режиме «Измеренный» реализован алгоритм по срабатыванию ИО частоты, который пускается при превышении ИО частоты во всех трех фазах, наличии разрешения ЧАПВ при отсутствии блокировки по напряжению.

\item
В устройстве реализована блокировка от многократных включений определяемая уставкой SGF4. Если уставка введена в работу, то при наличии команды включения выключателя от ЧАПВ происходит сброс триггеров готовности и разрешения ЧАПВ. Информация о блокировке повторного действия ЧАПВ будет отображаться выходным сигналом «ЧАПВ: Блок. повт. ЧАПВ».

\item
Сброс готовности ЧАПВ формируется при наличии одного из следующих условий:
\begin{itemize}
    \item функции АЧР и ЧАПВ выведены из работы;
    \item активный дискретный сигнал блокировки ЧАПВ;
    \item отсутствие сигнала «КСВ: РФК»;
    \item активный сигнал отключения выключателя от РЗ;
    \item любой из активных сигналов: «КСВ: Блокировка вкл», «Внешний запрет ЧАПВ»;
    \item принудительный сброс от кнопки сброса блокировки ЧАПВ «Сброс блокир. ЧПАВ».
\end{itemize}
\item
При наличии пусковых условий функции запускается отсчет выдержки времени таймера Т2, сбрасывается триггер готовности ЧАПВ и при отсутствии блокировки включения формируется команда на включение выключателя.

\item
Ввод функции в работу на этапе параметрирования устройства осуществляется программным переключателем SGF1 «Ввод\_функции», для работы также необходимо введенное состояние функции «АЧР», что сигнализируется активным статусом сигнала «АЧР: Ввод». Информация о введенной функции отображается выходным сигналом «ЧАПВ: Ввод».

\item
Функция может быть выведена из работы оперативно путем активации сигнала «ОВ ЧАПВ», что характеризуется активным состоянием сигнала «ЧАПВ: Оперативный вывод» на выходе функции. 
\item
Алгоритм работы функции выполнен в соответствии с рисунком \ref{chapv:img1}. В таблице \ref{chapv:tbl1} приведены параметры, необходимые для настройки функции.

\vspace{1mm}
\begin{figure}[H]
\centering
\includegraphics[width=1\textwidth,height=1\textheight,keepaspectratio]{\fbpath/120. ЧАПВ/_latex/img/chapv.pdf}
\captionof{figure}{Функциональная схема логики ЧАПВ}\label{chapv:img1}
\end{figure}

\end{enumerate}
